% EDITION: citation-based shortening, 4 October 2026.
% Earlier audits refer to the source edition; this revision has not been independently audited.
\section{P04: Arbitrary-\texorpdfstring{$(n,q)$}{(n,q)} Fourier spectrum and normality}
\label{sec:P04}
The standing assumptions are $n\ge5$, $1\le q<n$ and
$a,k,c,g,h>0$. Use the canonical domain $\Domain=(-1,\infty)^n$,
indices in $\Zn$, predecessor shift $(Pz)_j=z_{j-1}$ and the
unitarily normalized, positive-exponent Fourier vectors from
Section~\ref{sec:notation}. No coprimality condition is imposed.
This is supporting, known Fourier machinery specialized to the
verified homogeneous Jacobian, not a novelty claim.

For any synchronous state $x^*\one\in\Domain$, set $s=1+x^*>0$.
Here $s$ is the denominator, whereas the scalar state variable below
is $r$. Import Proposition~\ref{prop:P03-homogeneous}, without
re-deriving its Jacobian:
\begin{align}
 J&:=DF(x^*\one)=dI+CP+GP^{-q}-HP^q,\label{eq:P04-input}\\
 d&=a-2kx^*-\frac{cx^*}{s^2}-\frac{g}{s}-hx^*,
 &C&=\frac{c}{s},\quad G=\frac{gx^*}{s^2},\quad H=hx^*.
 \label{eq:P04-coefficients}
\end{align}
All coefficients are real; $C>0$, while $G,H$ have the sign of
$x^*$ (and vanish at zero). The spectral statements through the
synchronous-mode identity apply even when the state is not an
equilibrium. At a nonstationary state they describe its Jacobian,
not an autonomous linearization about an equilibrium.

\subsection{Fourier diagonalization}
\label{slot:P04-1}

\begin{theorem}[Complete homogeneous Fourier spectrum]
\label{thm:P04-spectrum}
Let $\theta_m=2\pi m/n$, $v^{(m)}_j=n^{-1/2}e^{\ii\theta_m j}$,
and $Q:=\Fourier=[v^{(0)}\ \cdots\ v^{(n-1)}]$.
Then $Q^*Q=QQ^*=I$, and the full spectrum, listed with algebraic
multiplicity, is
\begin{equation}
 \boxed{\lambda_m=d+Ce^{-\ii\theta_m}
                      +Ge^{\ii q\theta_m}-He^{-\ii q\theta_m},
           \qquad m=0,\ldots,n-1.}
 \label{eq:P04-spectrum}
\end{equation}
In particular,
\begin{equation}
 Jv^{(m)}=\lambda_m v^{(m)},\qquad
 Q^*JQ=\Lambda:=\diag(\lambda_0,\ldots,\lambda_{n-1}),\qquad
 J=Q\Lambda Q^*.
 \label{eq:P04-diagonalization}
\end{equation}
These statements do not require distinct shift offsets, distinct
numerical eigenvalues, nonzero combined coefficients, or
$\gcd(n,q)=1$.
\end{theorem}
\begin{proof}
Equation~\eqref{eq:P04-input} is a polynomial in the cyclic shift,
so Gray's theorem \cite[Theorem~3.1]{Gray2006} gives a unitary Fourier
diagonalization, also when offsets or numerical eigenvalues coincide.
The project's convention fixes the phases as
\begin{equation}
 Pv^{(m)}=e^{-\ii\theta_m}v^{(m)},\quad
 P^{-q}v^{(m)}=e^{\ii q\theta_m}v^{(m)},\quad
 P^qv^{(m)}=e^{-\ii q\theta_m}v^{(m)}.
 \label{eq:P04-shifts}
\end{equation}
Substitution gives \eqref{eq:P04-spectrum} and
\eqref{eq:P04-diagonalization}, hence
\begin{equation}
 \det(\zeta I-J)=\prod_{m=0}^{n-1}(\zeta-\lambda_m).
 \label{eq:P04-characteristic}
\end{equation}
We reuse this diagonalization throughout.
\end{proof}

\subsection{Real and imaginary dispersion relations}

Writing $\lambda_m=r_m+\ii\omega_m$, Euler's identities give
\begin{align}
 \boxed{r_m=d+C\cos\theta_m+(G-H)\cos(q\theta_m),}
   \label{eq:P04-real}\\
 \boxed{\omega_m=-C\sin\theta_m+(G+H)\sin(q\theta_m).}
   \label{eq:P04-imag}
\end{align}
Indeed, the imaginary contributions of the three shift terms are
$-C\sin\theta_m$, $+G\sin(q\theta_m)$ and
$+H\sin(q\theta_m)$: the last plus sign is the product of
$-H$ and the negative imaginary part of $e^{-\ii q\theta_m}$.
Their real contributions are $C\cos\theta_m$,
$G\cos(q\theta_m)$ and $-H\cos(q\theta_m)$.
Here $\omega_m$ is a signed imaginary part, possibly zero; it is not
the reserved positive frequency $\omega$ of a future Hopf theorem.

\subsection{Conjugacy and multiplicities}

For a label $m$, write $\bar m=(-m)\bmod n$ in $\{0,\ldots,n-1\}$;
in particular $\bar0=0$, not an additional label $n$.
The phases at $\bar m$ are conjugates of those at $m$, so
\begin{equation}
 v^{(\bar m)}=\overline{v^{(m)}},\qquad
 \lambda_{\bar m}=\overline{\lambda_m},\qquad
 r_{\bar m}=r_m,\quad \omega_{\bar m}=-\omega_m.
 \label{eq:P04-conjugacy}
\end{equation}
This verifies conjugacy explicitly from the formula and real
coefficients, consistently with $J$ being real.

A label is self-conjugate precisely when $2m=0\pmod n$.
For odd $n$ there is the singleton $0$ and the $(n-1)/2$ pairs
$\{m,n-m\}$ with $1\le m\le(n-1)/2$.
For even $n$ the singletons are $0,n/2$, with $n/2-1$ pairs
$1\le m<n/2$. Each singleton necessarily has a real eigenvalue.
The converse is false: a non-singleton pair can have $\omega_m=0$,
and then its two labels give the same real eigenvalue.

\begin{proposition}[Numerical eigenvalue multiplicities]
\label{prop:P04-multiplicity}
For a numerical eigenvalue $\eta$, define its label set
$M_\eta=\{m:0\le m<n,\ \lambda_m=\eta\}$.
Its algebraic and geometric multiplicities over $\C$ both equal
$|M_\eta|$, and
\[
 \ker(J-\eta I)=\operatorname{span}_{\C}
                  \{v^{(m)}:m\in M_\eta\}.
\]
Every eigenvalue is semisimple. Each Fourier label contributes one
one-dimensional Fourier coordinate, but this does not imply that
its numerical eigenvalue is simple.
\end{proposition}
\begin{proof}
Equation~\eqref{eq:P04-characteristic} counts the algebraic
multiplicity. Write $z=Qw$; then $(J-\eta I)z=0$ if and only if
$(\lambda_m-\eta)w_m=0$ for every $m$. Exactly the coordinates in
$M_\eta$ are free, proving the eigenspace formula and dimension.
Diagonalization also excludes nontrivial Jordan blocks.
\end{proof}

Thus ``Fourier-mode multiplicity'' means the number of labels
producing a given value, if that term is used at all; it is not a
second source of multiplicity. A nonreal conjugate pair consists
of two different numerical eigenvalues, each of which can also
coincide with eigenvalues at other labels. A tie $r_m=r_\ell$
alone says nothing about equality of $\lambda_m,\lambda_\ell$.
Singleton labels can likewise share their numerical value with
other labels; singleton refers to label conjugacy, not simplicity.

For a concrete canonical example, take $q=1$, $x^*=1$ and
$(a,k,c,g,h)=(2,1,4,4,1)$. P02's scalar equation vanishes,
and P03 gives $d=-4,C=2,G=H=1$. Hence
$J=-4I+P+P^{-1}$ and $\lambda_m=-4+2\cos\theta_m$;
each non-singleton pair has a repeated real eigenvalue.
At the origin, by contrast, $G=H=0$, $C=c>0$ and $d=a-g$,
so \eqref{eq:P04-spectrum} reduces to P03's
$a-g+ce^{-\ii\theta_m}$. The distinct roots of unity and $c>0$
make all those origin eigenvalues simple, as established in
Proposition~\ref{prop:P03-origin-spectrum}.

\subsection{Offset collisions and self-conjugate modes}
\label{slot:P04-extra-1}

P03's exhaustive collision list requires no replacement spectrum.
Direct substitution in the master formula gives the following
checks, with phases understood modulo $2\pi$.
\begin{enumerate}[label=(\roman*),leftmargin=*]
\item If $q=1$, then
\[
 J=dI+(C-H)P+GP^{-1},\qquad
 \lambda_m=d+(C-H)e^{-\ii\theta_m}+Ge^{\ii\theta_m}.
\]
The same expression follows by combining the $C$ and $-H$ terms
in \eqref{eq:P04-spectrum}. Here
$r_m=d+(C+G-H)\cos\theta_m$ and
$\omega_m=(-C+G+H)\sin\theta_m$.
\item If $q=n-1$, then $P^{-q}=P$ and $P^q=P^{-1}$.
Since $e^{\ii(n-1)\theta_m}=e^{-\ii\theta_m}$ and
$e^{-\ii(n-1)\theta_m}=e^{\ii\theta_m}$,
\[
 J=dI+(C+G)P-HP^{-1},\qquad
 \lambda_m=d+(C+G)e^{-\ii\theta_m}-He^{\ii\theta_m}.
\]
Thus $r_m=d+(C+G-H)\cos\theta_m$ and
$\omega_m=-(C+G+H)\sin\theta_m$, also obtained from
\eqref{eq:P04-real}--\eqref{eq:P04-imag}.
\item If $n$ is even and $q=n/2$, then
$e^{\pm\ii q\theta_m}=e^{\pm\ii\pi m}=(-1)^m$ and
$P^{-q}=P^q=P^{n/2}$. Therefore
\[
 J=dI+CP+(G-H)P^{n/2},\qquad
 \lambda_m=d+Ce^{-\ii\theta_m}+(G-H)(-1)^m.
\]
Here $r_m=d+C\cos\theta_m+(G-H)(-1)^m$ and
$\omega_m=-C\sin\theta_m$.
\end{enumerate}
These equalities include cancellations such as $C=H$ in the first
case or $G=H$ in the third. Adding contributions before or after
Fourier transformation gives identical results; coincident entries
must never be overwritten.

The singleton formulas, for every admissible $q$, are
\begin{align}
 \lambda_0&=d+C+G-H\in\R,\label{eq:P04-zero-mode}\\
 \lambda_{n/2}&=d-C+(G-H)(-1)^q\in\R
       \quad(n\text{ even}).\label{eq:P04-half-mode}
\end{align}
For the second formula use $e^{-\ii\pi}=-1$ and
$e^{\pm\ii q\pi}=(-1)^q$; all sine terms vanish.
It equals $d-C+G-H$ for even $q$ and $d-C-G+H$ for odd $q$.
When also $q=n/2$, the factor is $(-1)^{n/2}$, agreeing with the
half-offset formula at $m=n/2$. Thus no real singleton critical
mode is omitted from subsequent all-mode tests.

\subsection{Normality}
\label{slot:P04-2}

\begin{proposition}[Normality]
\label{prop:P04-normality}
Every homogeneous Jacobian \eqref{eq:P04-input} is normal over
$\C^n$: $JJ^*=J^*J$.
\end{proposition}
\begin{proof}
The unitary diagonalization in Theorem~\ref{thm:P04-spectrum}
makes $J$ normal, as in \cite[\S3.1]{Gray2006}.
\end{proof}

Normality here concerns homogeneous states, including negative ones in the domain; it is not asserted for heterogeneous or nonsynchronous Jacobians.

\subsection{Synchronous mode}

P02's scalar field, using a fresh state variable $r$, is
\[
 f(r)=ar-(k+h)r^2+\frac{(c-g)r}{1+r},\qquad r>-1.
\]

Differentiate $F(r\one)=f(r)\one$ and use $v^{(0)}=\one/\sqrt n$:
\begin{equation}
 \boxed{\lambda_0=d+C+G-H=f'(x^*).}
 \label{eq:P04-synchronous-derivative}
\end{equation}
This holds at any synchronous state, without the equilibrium hypothesis.

\paragraph{Positive synchronous equilibrium specialization.}
Now assume additionally $x^*>0$ and $F(x^*\one)=0$.
Then $C,G,H>0$. Put $A=k+h$ as in the canonical notation.
P02 and P03 give
\[
 a-Ax^*+\frac{c-g}{s}=0,\qquad
 d=-kx^*-\frac{c}{s}-\frac{cx^*}{s^2}
   =-kx^*-\frac{c(1+2x^*)}{s^2}.
\]
This is an optional equilibrium-only replacement for $d$ in
\eqref{eq:P04-spectrum}; the unsimplified expression
\eqref{eq:P04-coefficients} remains valid at all synchronous states.
Writing $f(r)=rB(r)$ with
$B(r)=a-Ar+(c-g)/(1+r)$ and $B(x^*)=0$ yields
\begin{equation}
 \boxed{\lambda_0=-x^*\left(A+\frac{c-g}{s^2}\right)
          =\frac{x^*}{s}\bigl(a-A-2Ax^*\bigr).}
 \label{eq:P04-equilibrium-zero-mode}
\end{equation}
The first equality uses $f'(x^*)=x^*B'(x^*)$; for the second,
substitute $c-g=(Ax^*-a)s$.
The equilibrium identities require a nonzero equilibrium; they also
hold at a negative nonzero equilibrium in $\Domain$, but $G,H>0$
is then false. At the origin use \eqref{eq:P04-coefficients}, which
gives $\lambda_0=a+c-g$. No comparison with other $r_m$ or selection
of a dominant mode is made here.

\subsection{Linear stability and hyperbolicity}
\label{slot:P04-3}

In this subsection $x^*\one$ is a synchronous \emph{equilibrium}
in $\Domain$. Adopt the project's spectral-abscissa notation
\begin{equation}
 \sa(J)=\max_{0\le m<n}r_m.
 \label{eq:P04-abscissa}
\end{equation}
For $\dot z=Jz$, exponential asymptotic stability means there exist
$M\ge1$ and $\beta>0$ such that
$\norm{z(t)}_2\le Me^{-\beta t}\norm{z(0)}_2$ for all $t\ge0$.

\begin{corollary}[Exact linear stability criterion]
\label{cor:P04-stability}
The linearized system is exponentially asymptotically stable if and
only if
\begin{equation}
 \boxed{\sa(J)<0\quad\Longleftrightarrow\quad
   d+C\cos\theta_m+(G-H)\cos(q\theta_m)<0
      \quad\text{for every }m.}
 \label{eq:P04-stability}
\end{equation}
For all $t\ge0$, its exact induced Euclidean propagator norm is
\begin{equation}
 \boxed{\norm{e^{tJ}}_2=e^{t\sa(J)}.}
 \label{eq:P04-semigroup}
\end{equation}
If $\sa(J)=0$, the linear system is Lyapunov stable but not
asymptotically stable. If $\sa(J)>0$, it is unstable.
\end{corollary}
\begin{proof}
Theorem~\ref{thm:P04-spectrum} gives
$e^{tJ}=Q\diag(e^{t\lambda_m})Q^*$, whose induced Euclidean norm
is $\max_m|e^{t\lambda_m}|=e^{t\sa(J)}$ for $t\ge0$.
This proves the strict criterion. At zero abscissa a maximizing real
Fourier block preserves nonzero norm; at positive abscissa it grows
exponentially. Real and complex induced norms agree for a real matrix,
so these conclusions apply to the real system as well.
\end{proof}

\begin{proposition}[Hyperbolicity]
\label{prop:P04-hyperbolicity}
The equilibrium is hyperbolic if and only if $r_m\ne0$ for every
Fourier label. It is nonhyperbolic whenever any $r_m=0$, including
when $\omega_m\ne0$. Invertibility requires instead
$\lambda_m\ne0$ for every label, and is strictly weaker than
hyperbolicity.
\end{proposition}
\begin{proof}
The complete spectrum is \eqref{eq:P04-spectrum}; substitute it in
the definitions of hyperbolicity and invertibility. A nonzero purely
imaginary eigenvalue meets the invertibility test but fails the
hyperbolicity test. For an instance within the canonical family,
P03's origin example $n=6$, $a=2,c=1,g=5/2$ (any positive $k,h$,
any canonical $q$) has $r_1=r_5=0$ and nonzero imaginary parts,
whereas no eigenvalue is zero: $ce^{-\ii\theta_m}=1/2$ is impossible
by modulus. Thus the weaker criterion is strictly weaker here as well.
\end{proof}
Hyperbolicity alone is not stability: hyperbolic equilibria may have
positive real parts. Nonhyperbolicity and instability may coexist
when one mode has zero real part and another has positive real part.

\paragraph{Nonlinear implications and their limits.}
The following conclusions concern the nonlinear equilibrium in the
open domain $\Domain$, using P01's $C^\infty$ smoothness there:
\begin{enumerate}[label=(\roman*),leftmargin=*]
\item If $\sa(J)<0$, the equilibrium is locally asymptotically stable;
in fact it is locally exponentially stable.
\item If $\sa(J)>0$, it is unstable, even if some other eigenvalues
have zero real part.
\item If $\sa(J)=0$, linearization alone is generally inconclusive
about nonlinear stability or attraction. Normality only settles the
linear system's boundedness in this case.
\end{enumerate}

These are the standard linearized stability and instability results
\cite[Ch.~1]{GuckenheimerHolmes1983}: P01 supplies smoothness on the
open domain, and the complete spectrum verifies the respective sign
condition. The instability theorem allows other eigenvalues on the
imaginary axis. At zero abscissa, a normal linearization alone does not
decide nonlinear stability (compare $\dot y=-y^3$ and $\dot y=y^3$).

All stability implications above are local. They establish neither
global attraction nor a global basin, even given P01's separate
boundedness result. At a positive equilibrium one may choose the
local ball inside $\PosOrth$, so the local assertions apply there
as well. At the origin the conclusions above are stated in
$\Domain$; instability relative to a restricted state set is a
separate question unless its admissible perturbations are checked.
A Hurwitz Jacobian is sufficient for nonlinear local asymptotic
stability, not a necessary condition in general. Thus the rejected
blanket nonlinear equivalence C040 and global assertion C153 are
not reinstated.

\paragraph{Dominant modes and first loss of stability.}
At a fixed parameter point and specified synchronous equilibrium,
the dominant-label set is $\argmax_{0\le m<n}r_m$, retaining all
ties. This is the fixed-point ``selected-mode set'' reserved in
Section~\ref{sec:notation}; it need not contain an unstable mode.
A \emph{first unstable mode along a path} additionally needs a
specified equilibrium branch $x^*(\nu)\one$, a specified oriented
parameter path $\param(\nu)$, a prior interval of stability, and a
loss-of-stability crossing event. For a continuous branch/path a
finite boundary event from $\sa(J)<0$ must have $\sa(J)=0$;
calling it an instability crossing also requires positive real part
on the unstable side, not merely tangency at zero. Several labels
or pairs can participate, so uniqueness and transverse crossing
cannot be presumed. No such path, crossing, or $q=2$ selection
analysis is supplied by this module.

\subsection{Arbitrary offset and the step-\texorpdfstring{$q$}{q} subgraph}
\label{prop:P04-step-cycles}
Let $\kappa=\gcd(n,q)$, $n=\kappa n'$ and $q=\kappa q'$,
where $\gcd(n',q')=1$. Repeated addition of $q$ returns an index
to itself after the least positive $L$ for which $n\mid Lq$.
This is equivalent to $n'\mid Lq'$, hence to $n'\mid L$,
so $L=n/\kappa$. The step-$q$ permutation therefore has exactly
$\kappa$ disjoint cycles, each of length $n/\kappa$; they are the
residue classes modulo $\kappa$. Reversing the step gives the same
partition. Thus coprimality is required only to call this step-$q$
subgraph a single cycle. The full canonical coupling also includes
the step-one predecessor cycle, and the step-$q$ partition is not
an assertion of full-network disconnection or of a block decomposition
of $J$ by those residue classes.

The Fourier sampling grid and label count depend on $n$; parity of
$n$ determines whether the $n/2$ singleton exists. The general phases
and possible eigenvalue coincidences depend on $q$ modulo $n$ as
well as $n$ and the coefficients. The value of $\lambda_0$ is
independent of $n,q$ at fixed coefficients and state; the half-mode
value depends on the parity of $q$. No proof of diagonalization,
conjugacy, multiplicity, normality, or the all-mode stability and
hyperbolicity criteria requires coprimality. The actual spectral
abscissa generally depends on the grid and phases; P03's
$n,q$-independent origin threshold is a special case.
